% ORACLE A analytic audit for CODEX_SPEC_oracleA.md.
% This is a fragment intended to be read with proof.tex.  The exact finite
% checks are in oracleA_verify.py and the numerical reconnaissance is recorded
% in oracleA_exploration.md.  No numerical observation is used below as proof.

\subsection{The exact short--arc transform and the remaining bilinear form}
\label{ssec:oracleA-result}

Write $e(t)=e^{2\pi i t}$ and $\mathbb T=\mathbb R/\mathbb Z$.  For a
prime block $P<p\leq2P$ put
\[
 c_{p,a}=\frac{(\log p)F_p(a)}p,
 \qquad
 Q_P=\sum_{P<p\leq2P}\sum_{a=1}^{p-1}|c_{p,a}|^2,
\]
and, for an integer $M\geq2$, define the cyclic short--arc sum
\[
 A_{P,M}(x)=
 \sum_{\substack{P<p\leq2P,\ 1\leq a<p\\
             a/p\in(x,x+1/M]_{\mathbb T}}}c_{p,a}.
\]
Endpoints are immaterial.  Finally set
\[
 D_P(k)=\sum_{P<p\leq2P}\sum_{a=1}^{p-1}
          c_{p,a}e(-ka/p),\qquad k\in\mathbb Z.
\]

\begin{proposition}[Exact triangular-kernel identity]
\label{prop:oracleA-triangle}
With $\|t\|$ denoting distance to the nearest integer, one has
\begin{align}
 \int_{\mathbb T}|A_{P,M}(x)|^2\,dx
 &=\sum_{p,a}\sum_{q,b}c_{p,a}\overline{c_{q,b}}
   \left(\frac1M-\left\|\frac ap-\frac bq\right\|\right)_+
 \label{eq:oracleA-triangle}\\
 &=\sum_{k\in\mathbb Z}\gamma_M(k)|D_P(k)|^2,
 \label{eq:oracleA-fejer}
\end{align}
where
\[
 \gamma_M(0)=\frac1{M^2},\qquad
 \gamma_M(k)=\left(\frac{\sin(\pi k/M)}{\pi k}\right)^2
 \quad(k\ne0).
\]
Moreover Fourier inversion on each $\mathbb F_p$ gives the exact identity
\begin{equation}
 D_P(k)=\sum_{P<p\leq2P}(\log p)
 \left(\mathbf1_{\mathcal Z_p}(k\bmod p)-\frac{Z(p)}p\right).
 \label{eq:oracleA-D-incidence}
\end{equation}
\end{proposition}

\begin{proof}
For a point $\lambda\in\mathbb T$, the set of $x$ for which
$\lambda\in(x,x+1/M]_{\mathbb T}$ is a cyclic interval of length $1/M$.
The intersection of the two intervals belonging to $\lambda$ and $\mu$
has measure $(M^{-1}-\|\lambda-\mu\|)_+$, proving
\eqref{eq:oracleA-triangle}.  The $k$-th Fourier coefficient of an interval
of length $1/M$ has squared modulus $\gamma_M(k)$; Parseval proves
\eqref{eq:oracleA-fejer}.  Finally,
\[
 \frac1p\sum_{a=0}^{p-1}F_p(a)e_p(-ak)
 =\mathbf1_{\mathcal Z_p}(k),
\]
and removing $a=0$, for which $F_p(0)=Z(p)$, proves
\eqref{eq:oracleA-D-incidence}.
\end{proof}

The two notions of zero frequency are both visible here.  The modular
frequency $a=0$ has been removed and produces the vertical centre
$\sum_p(\log p)Z(p)/p$.  The integer Fourier mode $k=0$ remains in
\eqref{eq:oracleA-fejer}; since $b_0=1$, it equals
$D_P(0)=-\sum_p(\log p)Z(p)/p$ and is weighted by $M^{-2}$.
Thus it has neither been silently discarded nor identified with the
empirical mean on $I_N$.  If $\mu_{P,I}$ is that empirical mean, then
orthogonal projection gives
\[
 \sum_{k\in I_N}|B_P^{(w)}(k)-\mu_{P,I}|^2
 \leq\sum_{k\in I_N}|D_P(k)|^2,
\]
where $B_P^{(w)}(k)=\sum_{P<p\leq2P}(\log p)
\mathbf1_{\mathcal Z_p}(k\bmod p)$.

There is a small but consequential point in applying this identity.
At the scale printed in~\eqref{eq:short-arc},
$\gamma_N(N)=\gamma_N(2N)=0$, so that estimate at scale $N$ does not by
itself control the shell $I_N=(N,2N]$.  Uniformity in the scale repairs
the problem without a phase twist.

\begin{proposition}[Correct Gallagher reduction by dilation]
\label{prop:oracleA-gallagher}
Assume that~\eqref{eq:short-arc}, with a factor $M^{o(1)}$, holds for
every pair $(P,M)$ satisfying $\sqrt M<P\leq M$.  Then, uniformly for
$\sqrt N<P\leq N$,
\begin{equation}
 \sum_{N<k\leq2N}|D_P(k)|^2\ll N^{1+o(1)}Q_P.
 \label{eq:oracleA-weighted-dispersion}
\end{equation}
Consequently, together with the unconditional $Z(p)\ll p^{2/3}$, the
short--arc target implies $W(n)=o(n)$ for every $n$ and hence
$G_n=e^{o(n)}$.
\end{proposition}

\begin{proof}
First suppose $P>2\sqrt N$.  Apply~\eqref{eq:short-arc} with the scale
$M=4N$.  This is legitimate because $\sqrt{4N}<P\leq4N$.  For
$N<k\leq2N$,
\[
 \gamma_{4N}(k)
 =\left(\frac{\sin(\pi k/(4N))}{\pi k}\right)^2
 \geq\frac1{8\pi^2N^2}.
\]
Equations~\eqref{eq:oracleA-fejer} and~\eqref{eq:short-arc} therefore give
\[
 \sum_{N<k\leq2N}|D_P(k)|^2
 \leq8\pi^2N^2\int_{\mathbb T}|A_{P,4N}(x)|^2\,dx
 \ll N^{1+o(1)}Q_P.
\]
For $\sqrt N<P\leq2\sqrt N$, the ordinary additive large sieve, applied
on any interval of $N$ consecutive integers, gives
\[
 \sum_{N<k\leq2N}|D_P(k)|^2
 \leq (N+4P^2)Q_P\ll NQ_P.
\]
This proves~\eqref{eq:oracleA-weighted-dispersion}.

Finite Fourier Parseval gives
\begin{equation}
 Q_P=\sum_{P<p\leq2P}(\log p)^2
 \left(\frac{Z(p)}p-\frac{Z(p)^2}{p^2}\right)
 \leq\sum_{P<p\leq2P}\frac{(\log p)^2Z(p)}p.
 \label{eq:oracleA-Q-parseval}
\end{equation}
Thus $Q_P\ll P^{2/3}\log P$, while
\[
 \mu_P:=\sum_{P<p\leq2P}\frac{(\log p)Z(p)}p\ll P^{2/3}.
\]
By~\eqref{eq:oracleA-D-incidence}, the logarithmically weighted bad-prime
sum in this block is at most
\[
 \mu_P+|D_P(k)|
 \ll P^{2/3}+N^{1/2+o(1)}P^{1/3}
 \ll N^{5/6+o(1)}
\]
for every $k\in I_N$.  Summing the $O(\log N)$ dyadic blocks and using the
small-prime and companion-block estimates already invoked in
Proposition~\ref{prop:ap-bdh-suffices} proves the claimed all-index bound.
\end{proof}

\begin{remark}[What this does and does not imply]
\label{rem:oracleA-weighted-not-amtd}
The direct assertion in Subsection~\ref{ssec:amtd} that
\eqref{eq:short-arc} ``implies AMTD block by block'' is not the conclusion
of the displayed Fourier calculation.  Its right-hand side is
\eqref{eq:oracleA-Q-parseval}; it is not
$S(P,N)=\sum_p|\Omega_p|$, and its coefficients carry the weights
$\log p$.  In particular, when $p>N$, a residue in $\mathcal Z_p$ need
not have a representative in $I_N$.  Proposition~\ref{prop:oracleA-gallagher}
is the valid conclusion: a weighted dispersion estimate with the vertical
majorant $NQ_P$, which is already strong enough for the conclusion of
Proposition~\ref{prop:ap-bdh-suffices}.  No unproved comparison with
$S(P,N)$ is needed.
\end{remark}

\subsection{Palindromic phases and top-block clipping}

\begin{lemma}[Exact doublet Fourier phase]
\label{lem:oracleA-doublet-phase}
Suppose
$\mathcal Z_p=\{r,p-1-r\}$, choose $0\leq r\leq(p-1)/2$, and put
$h=p-1-2r$.  Then, for every $a\bmod p$,
\begin{equation}
 F_p(a)=2e_p\!\left(\frac{a(p-1)}2\right)
             \cos\!\left(\frac{\pi ah}{p}\right).
 \label{eq:oracleA-doublet-F}
\end{equation}
Here $e_p(a/2)$ in the palindromic reality statement means multiplication
by the inverse of $2$ in $\mathbb F_p$; consequently
$e_p(a/2)F_p(a)=2\cos(\pi ah/p)\in\mathbb R$.
\end{lemma}

\begin{proof}
The two integers $r$ and $p-1-r$ have average $(p-1)/2$ and difference
$h$.  Factoring their two exponentials proves~\eqref{eq:oracleA-doublet-F}.
Since $(p+1)/2$ represents $2^{-1}\pmod p$, multiplying the phase in
\eqref{eq:oracleA-doublet-F} by $e_p(a(p+1)/2)$ cancels it.
\end{proof}

In ordinary real arguments this real rotation is
$e(-a(p-1)/(2p))F_p(a)$, because $(p+1)/2\equiv-(p-1)/2\pmod p$.
The expression with the opposite ordinary sign is not generally real;
the distinction is checked on the actual doublet
$\mathcal Z_{17}=\{3,13\}$ by \texttt{oracleA\_verify.py}.

The phase $e(\theta(3p-1)/2)$ in
Remark~\ref{rem:palindromic-fourier} is a phase of two particular
\emph{integer lifts}; it is not the ordinary real phase used to make
$F_p(a)$ real.  At the discrete value $\theta=a/p$ it does coincide
modulo $1$ with the phase $e_p(a(p-1)/2)$ in
\eqref{eq:oracleA-doublet-F}, since the two centres differ by $p$.
For continuous $\theta$, however, it is the centre of the chosen integer
lifts and therefore changes when those lifts change.  It also has a
restricted validity domain.  In general
the contribution of a residue $u\bmod p$ to the shell is the finite set
\begin{equation}
 H_{N,p}(u)=\{u+\ell p:N<u+\ell p\leq2N,\ \ell\in\mathbb Z\}.
 \label{eq:oracleA-hit-lifts}
\end{equation}
The shell Fourier sum must use
$H_{N,p}(r)\cup H_{N,p}(p-1-r)$, including every admissible lift.
For example, when $p\in(N,2N]$, the two displayed lifts
\[
 m_1=p+r,\qquad m_2=2p-1-r
\]
both belong to $I_N$ precisely when
\begin{equation}
 2p-1-2N\leq r\leq2N-p.
 \label{eq:oracleA-two-lift-domain}
\end{equation}
Only on this subfamily does
\[
 e(\theta m_1)+e(\theta m_2)
 =2e\!\left(\theta\frac{3p-1}{2}\right)
       \cos(\pi\theta h)
\]
hold for the contribution of the actual hit set $\Omega_p\cap I_N$.
Outside~\eqref{eq:oracleA-two-lift-domain}, one or both lifts are clipped;
when the reflected residue $p-1-r$ itself lies in $I_N$, the relevant pair
has centre $(2p-1)/2$, not $(3p-1)/2$.  Parseval still gives
\[
 \int_{\mathbb T}\left|\sum_p\sum_{m\in\Omega_p\cap I_N}e(\theta m)
 \right|^2d\theta
 =T+F,
\]
but only when the sum uses the actual, clipped hit positions.  Thus the
unqualified formula in Remark~\ref{rem:palindromic-fourier} cannot be the
starting point for all top-block doublets.  In the computational block
$N/2<p\leq N$, there can additionally be quotient-$2$ or quotient-$3$
lifts in~\eqref{eq:oracleA-hit-lifts}; these are also absent from the
complete two-term sum in that remark.

\subsection{The exact cross-prime lemma and a sharp obstruction}

Put $K_M(t)=(M^{-1}-\|t\|)_+$.  Splitting
\eqref{eq:oracleA-triangle} into equal and unequal primes gives
\[
 \int_{\mathbb T}|A_{P,M}|^2
 =\frac{Q_P}{M}+R_{=}(P,M)+R_{\ne}(P,M),
\]
where
\begin{equation}
 R_{\ne}(P,M)=
 \sum_{\substack{P<p,q\leq2P\\p\ne q}}
 \sum_{a=1}^{p-1}\sum_{b=1}^{q-1}
 K_M\!\left(\frac ap-\frac bq\right)
 \frac{(\log p)(\log q)}{pq}
 F_p(a)\overline{F_q(b)}.
 \label{eq:oracleA-cross-bilinear}
\end{equation}
Since $p\leq2P\leq2M$, each $a/p$ has at most two same-prime neighbours
in an arc of length $1/M$, and $|R_{=}(P,M)|\leq2Q_P/M$ by
$2|uv|\leq|u|^2+|v|^2$.  Hence the precise missing analytic lemma is the
one-sided scalar estimate
\begin{equation}
 \operatorname{Re}R_{\ne}(P,M)\leq C M^{-1+o(1)}Q_P.
 \tag{SDC}\label{eq:oracleA-sdc}
\end{equation}
Equivalently, this asks for cancellation among the near-Farey solutions
\[
 \left\|\frac ap-\frac bq\right\|<\frac1M,
 \qquad
 \min\{|aq-bp|,\,pq-|aq-bp|\}<\frac{pq}{M}.
\]
This is a scalar Rayleigh bound for the actual Ap\'ery vector; it neither
asks for a full operator norm nor discards the signs as Cauchy--Schwarz
does.  It is the exact bilinear input still missing from the Linnik route
of Remark~\ref{rem:linnik}.

There is also an exact reciprocal-fraction parametrization.  Let
$\langle d\rangle_{pq}$ be the representative of $d\bmod pq$ of least
absolute value, and let $\bar q$ and $\bar p$ denote inverses modulo $p$
and $q$, respectively.  For $p\ne q$, putting
$k=\langle aq-bp\rangle_{pq}$ is a bijection between the pairs occurring
in~\eqref{eq:oracleA-cross-bilinear} and
\[
 0<|k|<\frac{pq}{M},\qquad (k,pq)=1,
 \qquad a\equiv k\bar q\pmod p,\quad
 b\equiv-k\bar p\pmod q.
\]
Since $F_q(-u)=\overline{F_q(u)}$, one obtains
\begin{align}
 R_{\ne}(P,M)&=\frac1M\,\mathcal B_{P,M},
 \label{eq:oracleA-reciprocal-form}\\
 \mathcal B_{P,M}
 &=\sum_{\substack{P<p,q\leq2P\\p\ne q}}
   \frac{(\log p)(\log q)}{pq}
   \sum_{\substack{0<|k|<pq/M\\(k,pq)=1}}
   \left(1-\frac{M|k|}{pq}\right)
   F_p(k\bar q)F_q(k\bar p).
 \label{eq:oracleA-moving-reciprocal}
\end{align}
Thus~\eqref{eq:oracleA-sdc} is precisely the one-sided moving
reciprocal-fraction estimate
$\operatorname{Re}\mathcal B_{P,M}\leq C M^{o(1)}Q_P$.
The numerators $k\bar q$, $k\bar p$ and both zero sets move jointly with
$p,q$; palindromy only replaces each doublet transform by a real cosine
after a moving phase rotation.  It does not turn
\eqref{eq:oracleA-moving-reciprocal} into a fixed-numerator bilinear form.

\begin{proposition}[One trace-function mismatch removed on doublets]
\label{prop:oracleA-doublet-trace}
On the doublet sector, $a\mapsto F_p(a)$ is the trace function of a
rank-$2$, uniformly bounded-conductor sheaf: explicitly it is the direct
sum of the two rank-$1$ Artin--Schreier sheaves with traces
$a\mapsto e_p(ar_p)$ and
$a\mapsto e_p(-a(1+r_p))$.  Thus the growing-complexity obstruction for
representing the \emph{local Fourier coefficient} disappears completely
on this sector.
\end{proposition}

\begin{proof}
For $\mathcal Z_p=\{r_p,p-1-r_p\}$,
\[
 F_p(a)=e_p(ar_p)+e_p(-a(1+r_p)).
\]
Each summand is a linear additive-character trace of rank and conductor
bounded independently of both $p$ and the slope.  Direct sum proves the
claim.
\end{proof}

This does not make the frameworks discussed in Remark~\ref{rem:linnik}
applicable as stated.  The slope $r_p$ is selected by a same-characteristic
Ap\'ery zero condition and moves with $p$; there is no fixed sheaf over one
fixed finite field whose extension-degree traces give this horizontal
family.  Moreover~\eqref{eq:oracleA-moving-reciprocal} couples two varying
characteristics at once.  The anchored family below has exactly the same
rank-$2$ conductor bound and still violates the target.  Hence bounded
local conductor removes one of the four mismatches listed in
Remark~\ref{rem:linnik}, but horizontal control of the moving slopes
remains the load-bearing arithmetic input.

Palindromic reality and the bound $Z(p)=O(1)$ cannot prove
\eqref{eq:oracleA-sdc}.  The following construction quantifies the
obstruction at exactly the required short-arc scale.

\begin{proposition}[Reflected anchored obstruction]
\label{prop:oracleA-anchored}
Let
\[
 \mathcal P_N=\{p\text{ prime}:N<p\leq2N,\
                     p\nmid(4N+1)(2N+1)\}
\]
and, for $p\in\mathcal P_N$, set
\[
 \mathcal Z_p^*=\{\,2N\bmod p,\,-2N-1\bmod p\,\}.
\]
Then every $\mathcal Z_p^*$ is a two-element set invariant under
$r\mapsto p-1-r$; it contains neither $0$ nor adjacent residues.  Yet the
short--arc estimate at $(P,M)=(N,4N)$ fails by a factor
$N^{1-o(1)}$.  Simultaneously, all columns contain the point
$2N\in I_N$, so their centered block dispersion is $\gg N^2$ in the
logarithmically weighted normalization.
\end{proposition}

\begin{proof}
The divisors of $4N+1$ are exactly those for which the two displayed
residues coincide; excluding divisors of $2N+1$ also keeps $0$ out of
the sets and rules out one of the two possible adjacency congruences.
The other would force $p\mid2N$, impossible for a prime $p\in(N,2N)$.
These exclusions remove only $O(1)$ primes from this dyadic range.  For
the corresponding Fourier coefficients, let
$D_N^*$ denote~\eqref{eq:oracleA-D-incidence}.  Since
$2N\bmod p\in\mathcal Z_p^*$ for every $p\in\mathcal P_N$, the prime
number theorem gives
\[
 D_N^*(2N)=\sum_{p\in\mathcal P_N}(\log p)(1-2/p)\asymp N.
\]
At $M=4N$ the Fej\'er weight of this mode is
\[
 \gamma_{4N}(2N)=\frac1{4\pi^2N^2}.
\]
Therefore~\eqref{eq:oracleA-fejer} gives
\[
 \int_{\mathbb T}|A^*_{N,4N}(x)|^2dx\gg1.
\]
On the other hand, finite Fourier Parseval gives
\[
 Q_N^*=\sum_{p\in\mathcal P_N}(\log p)^2
          \left(\frac2p-\frac4{p^2}\right)\asymp\log N.
\]
Thus the proposed right-hand side is
$(4N)^{-1+o(1)}Q_N^*=N^{-1+o(1)}\log N=o(1)$.
The same common hit at $2N$ gives a centered deviation $\asymp N$ and
hence weighted dispersion $\gg N^2$: indeed every column has at most two
hits in the interval of length $N$, so the empirical weighted mean is
$O(N^{-1}\sum_{p\in\mathcal P_N}\log p)=O(1)$.
\end{proof}

The obstruction is not an assertion about the Ap\'ery zeros.  It proves
the required impossibility content: reflection, the exact palindromic
phase, doublet cardinality, vertical Parseval, and generic prime spacing
do not imply the short--arc bound.  A successful proof must insert an
Ap\'ery-specific horizontal fact into~\eqref{eq:oracleA-cross-bilinear},
exactly as anticipated in Remark~\ref{rem:linnik}; the separately running
arithmetic oracle is responsible for such an input.

\subsection{Computational reconnaissance}

The reproducible computation in \texttt{oracleA\_explore.py} and
\texttt{oracleA\_exploration.md} uses the exact zero-set bank and the
blocks $N/2<p\leq N$ for
$N=8192,16384,32768,65536$.  For the complete two-position doublet sum,
the exact Parseval ratios $(T+F)/T$ are between $1.0224$ and $1.0319$.
After declaring the major arcs
$|\theta-a/q|\leq16/(qN)$ with $q\leq16$, the largest sampled minor-arc
peaks are between $3.38\sqrt T$ and $3.69\sqrt T$.  The largest
nondegenerate connected correlation between
$e(3p\theta/2)$ and $\cos(\pi\theta h_p)$ at those small-denominator
rationals falls from $0.2204$ to $0.0657$ across the four scales.
An independent-uniform-gap collision model predicts the observed
complete-pair collision count within ratios $0.756$ to $1.128$.

The shell computation separately enumerates every lift in
\eqref{eq:oracleA-hit-lifts}.  With logarithmic weights its true
off-diagonal ratios $E^\circ/S_c$ are respectively
$0.00497,0.00185,-0.00689,0.00340$.  Direct counting and an independent
finite-group Fourier reconstruction agree within $4.4\cdot10^{-10}$.
The clipping distinction is numerically substantial: at $N=8192$ the
complete two-position sum retains only $162$ of its displayed positions
in $I_N$, whereas the actual doublet columns have $374$ shell lifts.
These observations support square-root-scale cancellation, but none is
used to infer~\eqref{eq:short-arc}.

\paragraph{STALL REPORT.}
The target~\eqref{eq:short-arc} is not proved, either for the full Ap\'ery
sum or for its doublet sub-sum.  The rigorous analytic outputs are:
\begin{enumerate}
 \item the exact triangular-kernel/Fej\'er identity
       \eqref{eq:oracleA-triangle}--\eqref{eq:oracleA-D-incidence};
 \item the corrected dilation argument proving that the stated uniform
       short--arc target really would imply the fully unconditional
       all-index conclusion, without the unsupported identification of
       its Fourier energy with $S(P,N)$;
 \item the exact doublet phase and the clipping domain
       \eqref{eq:oracleA-two-lift-domain}, correcting the unrestricted
       use of Remark~\ref{rem:palindromic-fourier};
 \item the one-sided cross-prime short-determinant estimate
       \eqref{eq:oracleA-sdc}, isolated as the precise missing bilinear
       lemma after the harmless same-prime terms are removed;
 \item Proposition~\ref{prop:oracleA-doublet-trace}, which removes the
       growing-conductor trace-function mismatch on the doublet sector
       while locating the surviving obstruction in the horizontal
       distribution of the moving slopes $r_p$;
 \item the reflected anchored family of
       Proposition~\ref{prop:oracleA-anchored}, which proves that none of
       palindromy, doublet dominance, vertical bounds, or generic large
       sieve geometry can supply that lemma.
\end{enumerate}
